% META: {"id":"TCOMPL-LOG-EX-002","chapitre":"TCOMPL-LOGARITHME-HISTORIQUE","type_objet":"exercice","capacites":["TCOMPL-LOG-C2"],"capacites_codes":["C2"],"methodes":[],"parcours":2,"competences":["calculer"],"duree_min":15,"mode_creation":"ex_nihilo","sources_inspiration":[],"fichier_tex":"chapitres/TCOMPL-LOGARITHME-HISTORIQUE/exercices/TCOMPL-LOG-EX-002.tex","corrige_tex":"chapitres/TCOMPL-LOGARITHME-HISTORIQUE/corriges/TCOMPL-LOG-CO-002.tex","status":"approved","origin":"generated","status_history":["generated","audited","accepted","approved"],"release_acceptance":"RELEASE_OWNER_BATCH_ACCEPTANCE","acceptance_closure_digest":"sha256:9b3ccf9a81c5520fbb7e03b7b2d3e7bf2057a02834b6d908bf3be5f49f3b553f","acceptance_receipt_digest":"sha256:4fad86f0282e0fa4e0419cca1ad6379ae7af5a280c04a4a90880f90a1c044242"}
% BEGIN-VERIFY
% from sympy import symbols, ln, solve, Eq
% x = symbols('x')
% sols = solve(Eq(ln(2*x+1), ln(x) + ln(3)), x)
% assert sols == [1]
% END-VERIFY
\begin{exercice}{TCOMPL-LOG-EX-002}{2}{15}
Résoudre l'équation $\ln(2x+1) = \ln(x) + \ln(3)$ (on précisera le domaine de définition).
\end{exercice}
